\documentclass[11pt,a4paper]{amsart}
\usepackage[T1]{fontenc}
\usepackage{amssymb}
\usepackage[hmargin=1in,vmargin=.9in]{geometry}
\usepackage[hidelinks]{hyperref}
\hypersetup{pdftitle={A counterexample to combination for finite-menu representations}}
\DeclareFontShape{OT1}{cmr}{bx}{n}{<-> cmbx10}{}
\DeclareMathSizes{8}{8}{8}{8}
\DeclareMathSizes{9}{9}{8}{8}
\DeclareMathSizes{10.95}{10.95}{8}{8}
\newtheorem{theorem}{Theorem}
\title{A counterexample to combination for finite-menu representations}
\date{}
\begin{document}
\begin{abstract}
Pointwise combination of two families of gamble cones need not
commute with their representation by desirable finite menus.
A two-outcome example uses a lexicographic cone and a sequence of
ordinary strict-expectation cones converging to it in membership.
Every cross-pair is inconsistent, while the two finite-menu assessments
have a common coherent extension. The discrepancy disappears for
closed membership families by a compactness corollary of the known
finitary conjunctive representation theorem.
\end{abstract}
\maketitle

\section{The combination question}

A gamble on a finite outcome space is a real function. A coherent
gamble cone \(D\) excludes zero, contains every nonzero nonnegative
gamble, and is closed under positive combinations.
For finite menus \(A\), put
\[
 K(D)=\{A:A\cap D\ne\varnothing\},\qquad
 \Gamma(\mathcal F)=\bigcap_{D\in\mathcal F}K(D).
\]
The inconsistent gamble cone is the whole gamble space and its
menu model contains all menus, including the empty one under the
null convention. Let \(\operatorname{cl}_D\) and
\(\operatorname{cl}_K\) denote the respective coherent natural
closures, with the null result when no coherent extension exists.
At a common label combine two families by
\[
 \mathcal F_1\otimes\mathcal F_2
 =\{\operatorname{cl}_D(D_1\cup D_2):
                         D_i\in\mathcal F_i,\ i=1,2\}.
\]
At different labels the individual gambles are extended to
their union label as functions constant in the new coordinates.

D\k{a}browska, Van Camp and Quaeghebeur
\cite[Prop.~11, \S4]{smps} prove the corresponding equality
with unrestricted menu sizes. Their conclusion asks about the
finite-menu case. We address the combination part, using the
specified original families and pointwise operation.

\begin{theorem}
There are nonempty families of coherent cones, one a singleton
and the other countably infinite, such that
\[
 \Gamma(\mathcal F_1\otimes\mathcal F_2)
 \ne
 \operatorname{cl}_K\bigl(\Gamma(\mathcal F_1)\cup
                          \Gamma(\mathcal F_2)\bigr).
\]
This holds on a common two-outcome domain, with a nonempty
singleton menu witnessing the inequality.
\end{theorem}

\section{Two outcomes and incompatible cones}

On outcomes \(0,1\) write \(f=(f_0,f_1)\) and
\(E_p(f)=(1-p)f_0+pf_1\).
Let \(p_n=1/2+1/(4n)\) for positive integers \(n\), and set
\[
 D_n=\{f:E_{p_n}(f)>0\}.
\]
Define \(L\) by lexicographic positivity of
\((E_{1/2}(f),f_1-f_0)\). In other words, \(f\in L\) if
\(E_{1/2}(f)>0\), or if that expectation is zero and \(f_1>f_0\).
All these cones exclude zero and include every nonzero nonnegative
gamble because the probability weights are positive. Positivity
in each rule is preserved by positive combinations. Use
\(\mathcal F_1=\{L\}\) and \(\mathcal F_2=\{D_n:n\ge1\}\).

For every fixed gamble,
\[
 E_{p_n}(f)=E_{1/2}(f)+\frac{f_1-f_0}{4n}.
\]
If its first term is nonzero, it determines the sign for large
\(n\). If it is zero, strict positivity holds exactly when the
second term is positive. Thus membership of each fixed gamble in
\(D_n\) eventually agrees with membership in \(L\).
If a finite menu avoids \(L\), take one \(n\) large enough for this
agreement for all of its elements. The menu then avoids \(D_n\)
and does not belong to \(\Gamma(\mathcal F_2)\).
Therefore
\(\Gamma(\mathcal F_2)\subseteq K(L)=\Gamma(\mathcal F_1)\).
The natural closure of the union is exactly the coherent \(K(L)\).

In the other order, for each \(n\) use
\(c_n=(1/2+p_n)/2\) and \(f_n=(c_n,c_n-1)\).
The two expectations are
\[
 E_{1/2}(f_n)=\frac1{8n},\qquad
 E_{p_n}(f_n)=-\frac1{8n}.
\]
So \(f_n\in L\) and \(-f_n\in D_n\).
Any cone containing both would contain the zero gamble by
addition. Every pairwise natural extension is thus inconsistent.
Its menu model contains the singleton \(\{(-1,-1)\}\), which
\(K(L)\) rejects. This proves the strict inequality.

\section{Finite menus and finite families}

The infinite menu \(\{-f_n:n\ge1\}\) hits every \(D_n\), since
its \(n\)-th element belongs to \(D_n\). It avoids \(L\), as
\(E_{1/2}(-f_n)<0\) for all \(n\). Allowing that menu breaks the
inclusion \(\Gamma(\mathcal F_2)\subseteq K(L)\). In the
infinitary closure, any choice \(-f_n\) from that menu conflicts
with the corresponding singleton \(\{f_n\}\) from \(K(L)\).
This is compatible with the infinitary Proposition~11.

Finiteness of a representing family is a different assumption
from finiteness of each menu. Earlier work describes how finite
representations can be combined to represent the least
compatible joint \cite[Props.~10--12]{marginal}. For the pointwise
equality here, the finite-family case follows directly from the
following compactness argument. Replacing an input family by
another representation can preserve its menu assessment without
preserving the original pairs of cones under combination.

\section{A compactness corollary of known representation theory}

We use the established finitary conjunctive representation
\cite[Cor.~37 and eq.~(24)]{rep}: the natural extension of a
finite-menu assessment is the intersection of all \(K(E)\)
for coherent \(E\) extending it as a menu assessment.
If no such \(E\) exists, this intersection is the null menu
model. The following application gives a sufficient condition
for the pointwise operation to commute.

View a cone as its indicator in the product
\(\{0,1\}^{\mathcal L}\), where \(\mathcal L\) is the gamble space.
Coherent cones form a closed subspace: the zero and positive-gamble
constraints are coordinate conditions, and each closure implication
for positive combinations involves finitely many coordinates.
Thus this subspace is compact. For finite \(A\), the condition
\(A\cap D\ne\varnothing\) is clopen. The closure \(\overline{\mathcal F}\)
of any nonempty family in this space has the same \(\Gamma\).

For a closed family \(\mathcal H\) of coherent cones, suppose
\(K(E)\) contains \(\Gamma(\mathcal H)\).
If every \(D\in\mathcal H\) has a gamble outside \(E\), the clopen
sets \(\{D:f\in D\}\) for \(f\notin E\) cover \(\mathcal H\).
A finite subcover gives a finite menu in \(\Gamma(\mathcal H)\)
missed by \(E\), a contradiction. Some member \(D\) is therefore
contained in \(E\). Conversely, \(D\subseteq E\) for a member \(D\)
implies \(K(E)\supseteq\Gamma(\mathcal H)\).
If labels differ, the argument applies to the restriction of \(E\)
to each label, under the same extension of gambles to the union.

Apply this equivalence to the closures of both input families.
Each combined cone model contains both input menu assessments.
Every coherent extension of the two menu assessments contains
the coherent natural extension of some cross-pair of cones.
Conversely, a coherent cross-pair extension hits every menu of
both input assessments. The coherent cones extending either
description are the same, after taking supersets. The conjunctive
representation theorem therefore yields
\[
 \Gamma(\overline{\mathcal F_1}\otimes\overline{\mathcal F_2})
  =\operatorname{cl}_K\bigl(\Gamma(\mathcal F_1)\cup
                            \Gamma(\mathcal F_2)\bigr).
\]
If all cross-pairs are inconsistent, both descriptions have no
coherent extension and both sides are null.
Finite families are closed in the product space.
In the example the closure adds the lexicographic limit \(L\)
and makes the cross-pair \((L,L)\) available.
This result concerns closing the families \emph{before}
combination. It follows from the prior representation theorem,
rather than constituting a separate homomorphism theorem for
the supplied arbitrary families.

\section{Verification and disclosure}

Internally reviewed, not peer reviewed. Exact rational replay
checks the pairwise conflict formula and explicit finite-menu
exclusions. The proof handles all real gambles and all menu
sizes. A partial Lean 4.30.0 certificate checks the algebraic
expectation identities, crossing signs, zero sum and rational
menu boundary. Coherence and compactness are proved in the text.

Under human direction, ChatGPT Astra 6.2 (OpenAI) generated the
discovery proof and its research packet. OpenAI Codex independently
reviewed the mathematics. GPT-6.2 Sol (OpenAI) checked the
literature and prepared the manuscript, exact replay, partial Lean certificate and
illustrations. DannyExperiments maintains the repository and
is the preferred citation identity. No AI system is proposed
as an author.

\begin{thebibliography}{9}
\raggedright
\bibitem{smps}
J. D\k{a}browska, A. Van Camp and E. Quaeghebeur,
\emph{Sets of desirable gamble sets and their representations as
information algebras}, SMPS (2026).
\href{https://arthurvancamp.github.io/proceeding/smps2026/JD-SMPS2026-paper.pdf}{Author paper}.
\bibitem{rep}
G. de Cooman, A. Van Camp and J. De Bock,
\emph{The logic behind desirable sets of things, and its
filter representation}, International Journal of Approximate
Reasoning \textbf{172} (2024), 109241.
\href{https://doi.org/10.1016/j.ijar.2024.109241}{doi:10.1016/j.ijar.2024.109241}.
\bibitem{marginal}
J. D\k{a}browska, A. Van Camp and E. Quaeghebeur,
\emph{The marginal problem for sets of desirable gambles sets
with a finite representation}, International Journal of
Approximate Reasoning \textbf{199} (2026), 109799.
\href{https://doi.org/10.1016/j.ijar.2026.109799}{doi:10.1016/j.ijar.2026.109799}.
\end{thebibliography}
\end{document}
